\section{Raw predictions throughout the mean-field-or-stronger regime}
\label{sec:raw-all-scales}

We now prove a finite-posterior oracle bound for the ordinary raw-output
posterior mean at every $\gamma\geq1$ and every eligible width, at any
fixed depth.  The overall prior scale decreases explicitly with sample
size.  Thus this theorem does not claim the same result at a fixed
prefactor $t=1$.

Use the model~\eqref{eq:effective-posterior}, with $H\geq2$, and the
effective energy $q_N$ from~\eqref{eq:effective-energy}.  Write
$a_N=\sigma\sqrt N$; the theorem below applies to any finite
$\sigma>0$ satisfying its stated bound on $a_N$.  To apply the preceding
lemmas, represent the same Gaussian law with $\gamma=1$ and $t=a_N$.
Recall the constants
$C_N,c_H$ from
\eqref{eq:effective-mgf-constant}.  Put
\begin{gather}
 T=64(1+c_H),\quad p=H/2,\quad J=\lceil\log_2(Tn)\rceil,
 \quad u_n=\frac1n\log\frac{2(J+1)}\delta,\notag\\
 D_n=\sqrt{2(d+1)\log(n+1)}+\frac{H-1}{2},\quad
 A_n=\frac{\sqrt2D_n\log^{3/2}(2n)}{H^p\sqrt n},\quad a_n=2^pA_n,
 \notag\\
 r_n=C_0[2\sqrt T(A_n+2\sqrt{u_n})+A_n^2+4u_n].
 \label{eq:raw-all-statistics}
\end{gather}
For each eligible $\lambda$, also define
\begin{equation}
 \nu_H=\frac{H-1}{2},\qquad b_H=\max\{1,\nu_H\},\qquad
 M_n=5T+r_n+\frac{512C_0^2Tu_n}{\lambda}.
 \label{eq:raw-all-risk-budget}
\end{equation}
Here $C_0\geq1$ is the same universal smooth-loss constant as in
\eqref{eq:finite-raw-sst}.  For a fixed teacher $g$ use $E,U,D_g,P_0,c_0$
from Lemma~\ref{lem:effective-teacher-tube}, and set
\[
 \rho=\frac\lambda{4(U+1)},\qquad
 \eta=\min\{1/2,\rho/(2P_0)\}.
\]
The fixed-teacher Bernstein term is
\[
 b_n(g,\delta)=(U+1)\sqrt{\frac{2R(g)\log(2/\delta)}n}
        +\frac{2(U+1)^2\log(2/\delta)}{3n}.
\]

\begin{theorem}[An all-width, all-exponent raw-output oracle bound]
\label{thm:raw-all-scales}
Let $H\geq2$, $n\geq2$, and $0<\delta<1$.  Suppose
\begin{gather}
 \frac1n\leq\lambda\leq1,\qquad
 \lambda^p\geq128C_0a_nT^{p-1/2},\qquad
 \lambda\geq512C_0u_n,
 \label{eq:raw-all-lambda}\\
 0<a_N\leq\min\left\{1,
       \frac{\lambda^{b_H}}
       {768\sqrt{M_n}(H-1)2^{\nu_H}T^{\max\{\nu_H-1,0\}}}\right\}.
 \label{eq:raw-all-prefactor}
\end{gather}
For every fixed width-$m$, weight-$B$ teacher $g$, with probability at
least $1-\delta$ over the sample, simultaneously for every $N\geq m$,
and every $\sigma>0$ satisfying~\eqref{eq:raw-all-prefactor},
\begin{align}
 \int R(z_w)\,\Pi(dw)
 &\leq R(g)+b_n(g,\delta)
  +\lambda\left[E+1+2a_N^2\left\{c_H
                 +D_g\log\frac1{c_0\eta}+\log2\right\}\right]
       \notag\\
 &\quad+r_n+\frac{512C_0^2Tu_n}{\lambda}
              +48\sqrt{M_n}a_N(H-1)2^{\nu_H}+66e^{-1/\lambda}.
 \label{eq:raw-all-oracle}
\end{align}
The same right side bounds the risk of the ordinary posterior mean.
It also bounds the usual clipped prediction conventions.
No limit in width or prior scale is taken.
\end{theorem}

For sufficiently large $n$, an explicit calibration is
\begin{gather}
 \lambda_n=\max\left\{1/n,
            [128C_0a_nT^{p-1/2}]^{1/p},512C_0u_n\right\},\qquad
 t_n=\frac{\lambda_n^{b_H}}
  {768\sqrt{M_n}(H-1)2^{\nu_H}T^{\max\{\nu_H-1,0\}}}.
 \label{eq:raw-all-schedule}
\end{gather}
For every $\gamma\geq1$ and every $N\geq m$, the prior
$\sigma^2=t_n^2N^{-\gamma}$ satisfies~\eqref{eq:raw-all-prefactor},
because $a_N=t_nN^{(1-\gamma)/2}\leq t_n$.  The same $t_n$ therefore
works simultaneously at all these widths and exponents.
At fixed $d,H,\delta$ these have orders
$\widetilde O(n^{-1/H})$ and $\widetilde O(n^{-b_H/H})$,
respectively; $M_n$ remains bounded as $n$ grows with these resources
fixed.  Thus the sufficient prefactor has exponent $1/2$ at $H=2$,
$1/3$ at $H=3$, and $(H-1)/(2H)$ for $H\geq4$.
The excess in~\eqref{eq:raw-all-oracle} is
$\widetilde O(n^{-1/H})$ for a fixed teacher.  The calibration depends
on $n,d,H,\delta$, not on the unknown data law or on a selected teacher.

\subsection{A weighted posterior tail for effective energy}

Let $F(w)=\widehat R_S(z_w)+\lambda q_N(w)$.  Every sample satisfies
\begin{equation}
 \int(1+q_N)\mathbf1_{\{F>T\}}\,d\Pi
       \leq33e^{-1/(\lambda\sigma^2)}\leq33e^{-1/\lambda}.
 \label{eq:raw-effective-tail}
\end{equation}
To prove this, the centered readout gives
$\mathbb E_\mu\widehat R_S\leq1+\sigma^2$ and hence
$Z\geq\exp[-(1+\sigma^2)/(2\lambda\sigma^2)]$.
Write $\beta=1/(2\lambda\sigma^2)$.  On $F>T$,
\[
 e^{-\beta\widehat R_S}
   \leq e^{-\beta T/16}e^{q_N/(32\sigma^2)},\qquad
 (1+q_N)e^{q_N/(32\sigma^2)}
    \leq(1+32\sigma^2)e^{q_N/(16\sigma^2)}
    \leq33e^{q_N/(16\sigma^2)}.
\]
The moment bound~\eqref{eq:effective-mgf} therefore makes the left
side of~\eqref{eq:raw-effective-tail} at most
\[
 33\exp\left(C_N+\frac{1+\sigma^2}{2\lambda\sigma^2}
                         -\frac{T}{32\lambda\sigma^2}\right).
\]
Since $\sigma^2C_N\leq c_H$, $\sigma^2\leq1$, and $\lambda\leq1$,
the exponent is at most
\[
 \frac{2c_H+2-T/16}{2\lambda\sigma^2}
       \leq-\frac1{\lambda\sigma^2}.
\]
This proves the claim.  Hard spectral truncation leaves the readout
unchanged, so both original and truncated outputs have absolute value
at most $\sqrt{q_N}$.  Thus the raw population risk on $F>T$ contributes
at most $66e^{-1/\lambda}$.

\subsection{Smooth-loss absorption on the low-objective region}

Apply~\eqref{eq:finite-raw-sst} to the deterministic energy classes
$\mathcal Z_{2^j}$, $0\leq j\leq J$, with failure probabilities
$\delta/[2(J+1)]$.  The common event has probability at least
$1-\delta/2$ and does not depend on $N$ or $\sigma$.
On $F\leq T$ write $q=q_N(w)$ and let $\widetilde w$ denote the
truncated network.  Then $q\leq T/\lambda\leq Tn$,
$\widehat R_S(z_w)\leq T$, and
Lemma~\ref{lem:effective-approximation} gives
\[
 e_w:=\|z_w-z_{\widetilde w}\|_\infty
       \leq8a_N(H-1)(1+q)^{\nu_H}\leq\frac{T}{96\sqrt{M_n}}
       \leq\sqrt T.
\]
For this use~\eqref{eq:raw-all-prefactor},
$1+T/\lambda\leq2T/\lambda$, $\lambda\leq1$, and $M_n\geq T$.
In particular, Minkowski yields
$\sqrt{\widehat R_S(z_{\widetilde w})}\leq2\sqrt T$.

For $q\leq1$, the unit-energy smooth-loss bound is at most $r_n$.
For $q>1$, a dyadic shell of radius at most $2q$ gives
\begin{align*}
 R(z_{\widetilde w})-\widehat R_S(z_{\widetilde w})
 \leq C_0[2\sqrt T a_nq^p+a_n^2q^{2p}
                  +4\sqrt{2T}\sqrt{qu_n}+8qu_n].
\end{align*}
The four terms are at most
\[
 \frac{\lambda q}{64},\quad
 \frac{\lambda q}{16384C_0},\quad
 \frac{\lambda q}{64}+\frac{512C_0^2Tu_n}{\lambda},\quad
 \frac{\lambda q}{64},
\]
respectively, by~\eqref{eq:raw-all-lambda}, $p\geq1$, and
$q\leq T/\lambda$.  Their sum is at most
$\lambda q/16+512C_0^2Tu_n/\lambda$.

For both energy cases, the preceding inequality implies
\[
 R(z_{\widetilde w})\leq
 \widehat R_S(z_{\widetilde w})+\lambda q/16
             +r_n+512C_0^2Tu_n/\lambda\leq M_n.
\]
Here $\widehat R_S(z_{\widetilde w})\leq4T$ and $\lambda q\leq T$.
Minkowski therefore transfers the two raw losses with cost at most
\begin{align*}
 R(z_w)-R(z_{\widetilde w})
  +\widehat R_S(z_{\widetilde w})-\widehat R_S(z_w)
 &\leq2e_w(\sqrt{M_n}+\sqrt T)+2e_w^2\\
 &\leq6\sqrt{M_n}e_w\\
 &\leq48\sqrt{M_n}a_N(H-1)(1+q)^{\nu_H}.
\end{align*}
If $q>1$ and $\nu_H\geq1$, use
$(1+q)^{\nu_H}\leq2^{\nu_H}(T/\lambda)^{\nu_H-1}q$;
if $0<\nu_H<1$, use $(1+q)^{\nu_H}\leq2^{\nu_H}q$.
In either case~\eqref{eq:raw-all-prefactor} bounds the transfer cost
by $\lambda q/16$.  For $q\leq1$ its bound is
$48\sqrt{M_n}a_N(H-1)2^{\nu_H}$.
Consequently, throughout $F\leq T$,
\begin{equation}
 R(z_w)\leq\widehat R_S(z_w)+\frac\lambda8 q_N(w)
       +r_n+\frac{512C_0^2Tu_n}{\lambda}
       +48\sqrt{M_n}a_N(H-1)2^{\nu_H}.
 \label{eq:raw-effective-deviation}
\end{equation}

\begin{proof}[Proof of Theorem~\ref{thm:raw-all-scales}]
Integrate~\eqref{eq:raw-effective-deviation} over $F\leq T$ and use
\eqref{eq:raw-effective-tail} outside.  The integrand
$\widehat R_S+\lambda q_N/8$ is nonnegative, so its restricted integral
is at most its full integral.  Apply the variational comparison
\eqref{eq:effective-variational} with the teacher-cylinder comparison
from~\eqref{eq:effective-comparison}.  The moment constant satisfies
$\sigma^2C_N\leq a_N^2c_H$.  These statements give the right side of
\eqref{eq:raw-all-oracle} with $\widehat R_S(g)$ in place of
$R(g)+b_n(g,\delta)$.

Bernstein's inequality for the fixed teacher supplies that replacement
except on a sample set of probability $\delta/2$.  The resulting event
is common to all the claimed widths and scales.  The posterior has
bounded density relative to its Gaussian prior, and
$|z_w|^2\leq Q_N(w)$; thus the ordinary posterior mean is integrable,
and Jensen proves its risk bound.  Clipping cannot increase squared
loss for labels in $[-1,1]$, giving the remaining conventions.
\end{proof}

The general scale condition also covers a fixed prefactor $t=1$
when $\gamma>1$ and
\[
 N\geq\left(\frac{768\sqrt{M_n}(H-1)2^{\nu_H}
                          T^{\max\{\nu_H-1,0\}}}
                         {\lambda_n^{b_H}}\right)^{2/(\gamma-1)}.
\]
At fixed $d,H,\delta,\gamma$, this sufficient width is
$\widetilde O(n^{2b_H/[H(\gamma-1)]})$.  It is an additional large-width
corollary, not an all-width fixed-prefactor theorem for
$1\leq\gamma<2$.

The coefficient of $R(g)$ is one.  A teacher may therefore be a
population-risk minimizer over any specified compact bounded-weight
class, chosen before sampling.  At fixed $d,H,\delta$, the theorem
competes with every such class whose width and weight caps are at most
$n^\alpha$, for every $\alpha<1/(4H)$: its excess tends to zero,
uniformly over every learner width above the class width and every
$\gamma\geq1$.  With a fixed weight cap the allowed range is
$\alpha<1/(2H)$.  This follows from $E=O(m^2B^2+B^2)$,
$U=mB$, the polynomial teacher-cylinder constants, and the fact that
the population minimizer has risk at most one.
